\chapter{Primitives and finite integration formulas}
\SourceRef{Zorich I, Sections 6.2--6.4, especially PDF pp. 377--386; Apostol, Section 7.20; Tao I, Section 11.9. For rationalizing substitutions compare Zorich I, Sections 5.7.4--5.7.5}
Unless stated otherwise $a<b$ and integration is the proper Riemann
integration of module 08. A scalar may be real or complex; order and exact
mean-value conclusions require real values.

\KEntry{K03.09.primitive}{Primitives are functions}
\begin{definition}
On an open interval $I$, a primitive of $f:I\to\K$ is a differentiable
function $F$ with $F'=f$ at every point. On a general interval impose
this on its interior, and specify continuity or one-sided derivatives at
endpoints separately. The indefinite integral denotes the family of all
primitives, not a number; the definite integral denotes a value under an
integrability hypothesis.
\end{definition}

\KEntry{K03.09.uniqueness}{Uniqueness up to a constant}
\begin{theorem}
Two primitives of the same function on an interval differ by a constant.
If endpoints are included, assume their continuous extensions there.
\end{theorem}
\begin{proof}
Their difference has derivative zero. For any two interior points, the
real mean value theorem applied between them gives equal values.
For complex values apply this separately to real and imaginary parts.
Continuity passes the equality to included endpoints. On a disconnected
domain the argument applies to each interval component, with independent
constants: for example a function equal to $0$ on $(-\infty,0)$ and $1$
on $(0,\infty)$ has zero derivative on its domain but is not constant there.
\end{proof}

\KEntry{K03.09.accumulation}{Lipschitz continuity of accumulation}
\begin{theorem}
If $f$ is bounded and integrable on $[a,b]$, then
$F(x)=\int_a^xf$ satisfies
$|F(y)-F(x)|\le M|y-x|$ for $M=\sup|f|$.
\end{theorem}
\begin{proof}
Restrictions are integrable by interval additivity. For $x\le y$,
$F(y)-F(x)=\int_x^yf$, whose modulus is at most $M(y-x)$ by
\KRef{K03.08.order}. Reversing $x,y$ proves the other case. This
argument uses bounded integrability and does not assume continuity of $f$.
\end{proof}

\KEntry{K03.09.ftc_first}{The first fundamental theorem}
\begin{theorem}
For integrable $f$ on $[a,b]$, its accumulation $F(x)=\int_a^xf$
is differentiable at every interior continuity point $x$ of $f$, with
$F'(x)=f(x)$. At a relatively continuous endpoint the corresponding
one-sided derivative equals that endpoint value.
\end{theorem}
\begin{proof}
For $h\ne0$ with $x+h\in[a,b]$, additivity and orientation give
\[\frac{F(x+h)-F(x)}h-f(x)=\frac1h\int_x^{x+h}(f(t)-f(x))\dd t.
\]
Its modulus is bounded by
$\sup_{t\text{ between }x\text{ and }x+h}|f(t)-f(x)|$, which tends to
zero by continuity at $x$. The same estimate with one-sided $h$ proves
the endpoint statements, with no two-sided derivative claimed there.
\end{proof}

\KEntry{K03.09.ftc_second}{Newton--Leibniz with integrable derivative}
\begin{theorem}
Suppose $F:[a,b]\to\K$ is continuous, differentiable on $(a,b)$,
and its derivative, extended by any finite endpoint values, is bounded
and Riemann integrable. Then
\[\int_a^bF'(x)\dd x=F(b)-F(a).
\]
\end{theorem}
\begin{proof}
For real $F$ and any partition the mean value theorem supplies
$\xi_i\in(x_{i-1},x_i)$ with
$F(x_i)-F(x_{i-1})=F'(\xi_i)\Delta x_i$. Summing telescopes to
$F(b)-F(a)$. As the mesh tends to zero, these sums tend to $\int F'$
by the all-tagged equivalence. For complex $F$ apply the real argument
to its components; their mean-value tags may differ, which causes no
problem because each real integral has the all-tagged limit.
Finite endpoint assignments do not change the integral.
\end{proof}

\KEntry{K03.09.continuous_primitive}{Existence for continuous functions}
\begin{theorem}
Every continuous scalar function on an interval $I$ has a primitive on
its interior, continuously defined at every included finite endpoint.
\end{theorem}
\begin{proof}
Fix $x_0\in I$ and define $F(x)=\int_{x_0}^xf$ on each compact
interval between $x_0$ and $x$. Continuity supplies proper integrability
there. Additivity ensures the definitions agree as the chosen compact
interval grows. The first fundamental theorem gives $F'=f$ at every
interior point. Local boundedness gives continuity at each included
endpoint by the accumulation estimate. No assertion of a finite value
at an excluded infinite or singular endpoint is involved.
\end{proof}

\KEntry{K03.09.parts}{Integration by parts}
\begin{theorem}
For $u,v\in C^1([a,b],\K)$,
\[\int_a^buv'=u(b)v(b)-u(a)v(a)-\int_a^bu'v.
\]
More generally it suffices that $u,v$ are continuous on $[a,b]$,
differentiable inside, and their derivatives admit bounded integrable
endpoint extensions.
\end{theorem}
\begin{proof}
The continuous $u,v$ are bounded integrable, so products with their
integrable derivatives are integrable by \KRef{K03.08.products}.
The product rule gives $(uv)'=u'v+uv'$ in the interior. The second
fundamental theorem applied to $uv$ and linearity give the formula.
The $C^1$ case satisfies all these hypotheses directly.
\end{proof}

\KEntry{K03.09.substitution}{Substitution for continuous integrands}
\begin{theorem}
If $\phi\in C^1([a,b],\R)$ and $f$ is continuous on an interval
containing $\phi([a,b])$, then
\[\int_a^bf(\phi(t))\phi'(t)\dd t
=\int_{\phi(a)}^{\phi(b)}f(x)\dd x.
\]
No monotonicity of $\phi$ is required in this version.
\end{theorem}
\begin{proof}
Let $m=\min\phi([a,b])$ and $M=\max\phi([a,b])$. Extend the restriction of $f$ to $[m,M]$ constantly beyond its endpoints: use $f(m)$ for $x<m$ and $f(M)$ for $x>M$. This gives a continuous scalar function $\widetilde f$ on $\R$, also when $m=M$. Choose its accumulation primitive $H$ on $\R$. This open-domain extension is needed even for a turning point such as $\phi(t)=t^2$ when the original $f$ was defined only on a closed interval.
The chain rule gives $(H\circ\phi)'=(\widetilde f\circ\phi)\phi'=(f\circ\phi)\phi'$, continuous
on $[a,b]$ with the one-sided endpoint convention. Newton--Leibniz gives
the left side as $H(\phi(b))-H(\phi(a))$. The same theorem for $H$
identifies this difference with the oriented integral on the right, because $\widetilde f=f$ between the two image endpoints.
\end{proof}

\KEntry{K03.09.monotone_substitution}{Substitution for Riemann-integrable integrands}
\begin{theorem}
If $\phi\in C^1([a,b],\R)$ is monotone and $f$ is bounded and Riemann
integrable on the interval between $\phi(a),\phi(b)$, then
$(f\circ\phi)\phi'$ is integrable and the same substitution formula holds.
\end{theorem}
\begin{proof}
First let $\phi$ be nondecreasing and nonconstant, $|f|\le M$, and
$L=\sup|\phi'|$. For a tagged partition $P$ put
$d_i=(\phi(t_i)-\phi(t_{i-1}))/\Delta t_i$.
The mean value theorem and uniform continuity of $\phi'$ imply
$|\phi'(\xi_i)-d_i|\le\omega_{\phi'}(|P|)$, where this uniform
modulus tends to zero. Thus the tagged sum for $g=(f\circ\phi)\phi'$
differs by at most $M(b-a)\omega_{\phi'}(|P|)$ from
\[\sum_i f(\phi(\xi_i))(\phi(t_i)-\phi(t_{i-1})).\]
After deleting repeated image points, this is a tagged sum for $f$ on
$[\phi(a),\phi(b)]$: terms on constant image cells are zero, and each
other image tag lies in its image cell. Its mesh is at most $L|P|$.
All such sums consequently approach $\int_{\phi(a)}^{\phi(b)}f$.
The displayed error tends to zero uniformly in tags, so $g$, which is
bounded, is integrable by \KRef{K03.08.tagged_equivalence} with the
required value. For constant $\phi$, $\phi'=0$ and both sides are zero.
For nonincreasing $\phi$, apply the increasing result to $-\phi$ and
$x\mapsto f(-x)$, then reverse orientation. Complex $f$ follows by
components. This proves existence of the transformed integral rather
than adding it as an extra assumption.
\end{proof}

\KEntry{K03.09.mean_value}{The weighted integral mean value theorem}
\begin{theorem}
For continuous real $f$ and integrable real $g\ge0$ on $[a,b]$, some
$c\in[a,b]$ satisfies $\int fg=f(c)\int g$.
\end{theorem}
\begin{proof}
Let $m=\min f$, $M=\max f$ and $G=\int g\ge0$.
Order gives $mG\le\int fg\le MG$. If $G=0$ these bounds make
$\int fg=0$, and any $c$ works. If $G>0$, the ratio $(\int fg)/G$
lies in $[m,M]$, and the intermediate value theorem supplies $c$.
The exact point conclusion is real; for complex $f$ the analogous
norm bounds hold but a point of equality need not exist.
\end{proof}

\KEntry{K03.09.second_mean_value}{The normalized second mean value theorem}
\begin{theorem}
If $f$ is real integrable and $g$ is nonnegative and nonincreasing on
$[a,b]$, then some $\xi\in[a,b]$ satisfies
\[\int_a^bg(x)f(x)\dd x=g(a)\int_a^\xi f(x)\dd x.
\]
\end{theorem}
\begin{proof}
The bounded monotone $g$ and its product with $f$ are integrable.
Let $F(x)=\int_a^xf$, continuous by the accumulation estimate, with
minimum $m$ and maximum $M$; $m\le F(a)=0\le M$.
For a partition put $g_i=g(x_i)$. The difference between $\int gf$ and
$\sum_{i=1}^ng_{i-1}(F(x_i)-F(x_{i-1}))$ is bounded by
\[\|f\|_\infty\sum_i\osc(g,[x_{i-1},x_i])\Delta x_i
\le\|f\|_\infty |P|(g(a)-g(b))\longrightarrow0.
\]
Finite summation by parts writes that sum as
\[g_{n-1}F(b)+\sum_{i=1}^{n-1}(g_{i-1}-g_i)F(x_i).
\]
All weights are nonnegative and sum to $g(a)$. It is consequently
between $mg(a)$ and $Mg(a)$, and the same holds for its limit $\int gf$.
If $g(a)=0$, monotonicity and nonnegativity give $g=0$. Otherwise
$(\int gf)/g(a)\in[m,M]$ and continuity of $F$ gives the required
$\xi$. The proof allows jumps and retains the actual endpoint value
$g(a)$, as in Zorich I, formula (6.33); it does not silently assume $C^1$.
\end{proof}

\KEntry{K03.09.elementary_primitives}{Elementary primitives with domains}
\begin{proposition}
The following primitive formulas hold on any interval in the indicated domain:
\begin{align*}
\int x^\alpha\dd x&=x^{\alpha+1}/(\alpha+1)+C
 &&(x>0,\ \alpha\ne-1),\\
\int x^{-1}\dd x&=\log|x|+C&&(x\ne0),\\
\int E(ax)\dd x&=E(ax)/a+C&&(a\ne0),\\
\int\cos x\dd x&=\sin x+C,&
\int\sin x\dd x&=-\cos x+C,\\
\int\frac{\dd x}{1+x^2}&=\arctan x+C,&
\int\frac{\dd x}{\sqrt{1-x^2}}&=\arcsin x+C\quad(|x|<1),\\
\int\tan x\dd x&=-\log|\cos x|+C&&(\cos x\ne0).
\end{align*}
For integer $\alpha\ge0$ the power formula holds on all of $\R$.
\end{proposition}
\begin{proof}
Differentiate each right side using module 07. In the logarithmic formula,
on $x<0$ the chain rule gives $(\log(-x))'=(-1)/(-x)=1/x$;
on $x>0$ the usual logarithm derivative applies. Likewise
$(\log|\cos x|)'=-\sin x/\cos x$ on each pole-free interval.
The other displayed derivatives are respectively $x^\alpha$, $E(ax)$,
$\cos x$, $\sin x$, $(1+x^2)^{-1}$ and $(1-x^2)^{-1/2}$.
Uniqueness of primitives gives the additive constants on each interval.
Constants on different sides of a pole are independent. The excluded
case $a=0$ is the constant integrand one, with primitive $x+C$.
\end{proof}

\KEntry{K03.09.trig_substitution}{The half-angle rationalization}
\begin{proposition}
On $I_k=((2k-1)\pi,(2k+1)\pi)$, put $t=\tan(x/2)$.
This is a $C^1$ increasing bijection onto $\R$, with inverse
$x=2\arctan t+2k\pi$, and
\[\sin x=\frac{2t}{1+t^2},\qquad
\cos x=\frac{1-t^2}{1+t^2},\qquad
\dd x=\frac{2\dd t}{1+t^2}.
\]
\end{proposition}
\begin{proof}
The tangent branch results give the bijection and its inverse.
For $u=x/2$, divide $\sin2u=2\sin u\cos u$ and
$\cos2u=\cos^2u-\sin^2u$ by
$1=\cos^2u(1+\tan^2u)$ to obtain the identities.
Differentiating the inverse gives $\dd x/\dd t=2/(1+t^2)$.
The proper substitution theorem applies on compact subintervals avoiding
the endpoints and all zeros of the original rational denominator.
Thus a rational expression $R(\sin x,\cos x)$ becomes rational in $t$,
but a primitive must be continued separately across the excluded branch
endpoints or denominator poles if continuation is possible.
\end{proof}

\KEntry{K03.09.algebraic_substitution}{Selected real square-root substitutions}
\begin{example}
For $A>0$, on $-A<x<A$, $x=A\sin u$ with
$u\in(-\pi/2,\pi/2)$ gives
$\sqrt{A^2-x^2}=A\cos u>0$ and $\dd x=A\cos u\dd u$.
On $x>A$, $x=A\cosh u$, $u>0$, gives
$\sqrt{x^2-A^2}=A\sinh u>0$ and $\dd x=A\sinh u\dd u$.
On all real $x$, $x=A\sinh u$ gives
$\sqrt{x^2+A^2}=A\cosh u>0$ and $\dd x=A\cosh u\dd u$.
\end{example}
\begin{proof}
Each radicand identity follows from the trigonometric or hyperbolic
square identity. The branch restrictions fix the sign of the square root
and make the substitution a bijection. The differentials follow by
differentiation. For a linear root $y=\sqrt{\alpha x+\beta}$ with
$\alpha\ne0$ and $\alpha x+\beta>0$, use
$x=(y^2-\beta)/\alpha$, $y>0$, and $\dd x=2y\dd y/\alpha$.
Completing the square, when the residual constant is nonzero, reduces
a quadratic with nonzero leading coefficient to a positive scaling of
one of the three displayed radicands on intervals where it is positive,
or reveals that no positive real branch exists. For a zero residual
constant, $\sqrt{A(x-h)^2}=\sqrt A\,|x-h|$ if $A>0$; on each side
of $h$ this is a fixed linear branch. If $A<0$ there is no open interval
of positive radicand. For instance
$\int\dd x/\sqrt{x^2+A^2}=u+C=\operatorname{arsinh}(x/A)+C$.
All reductions use the substitution theorem on regular compact
subintervals. These selected reductions do not assert a universal
algorithm for elementary antiderivatives.
\end{proof}

\KEntry{K03.09.holder}{H\"older's finite integral inequality}
\begin{theorem}
Let $p,q>1$, $1/p+1/q=1$, and let $f,g$ be integrable real or complex
functions on $[a,b]$. Then
\[\int_a^b|fg|\le
\left(\int_a^b|f|^p\right)^{1/p}
\left(\int_a^b|g|^q\right)^{1/q}.
\]
\end{theorem}
\begin{proof}
All absolute powers are integrable by continuous composition on bounded
ranges. Young's pointwise inequality $uv\le u^p/p+v^q/q$ for $u,v\ge0$
follows as follows: for fixed $v>0$ the derivative of
$u^p/p-uv+v^q/q$ is $u^{p-1}-v$, so its minimum is at
$u=v^{1/(p-1)}$ and has value zero because $q=p/(p-1)$.
The case $v=0$ is immediate.
Put $A=(\int|f|^p)^{1/p}$, $B=(\int|g|^q)^{1/q}$.
If $A,B>0$, apply Young to $|f|/A,|g|/B$ and integrate, obtaining
$\int|fg|/(AB)\le1/p+1/q=1$.
If $A=0$, Young applied to $|f|/\delta$ and $|g|$ gives
$\int|fg|\le\delta\int|g|^q/q$ for every $\delta>0$, so the left
side is zero. If $B=0$, interchange the roles. This handles zero
seminorms without claiming the functions vanish at every point.
The case $p=q=2$ is Cauchy--Schwarz.
\end{proof}

\KEntry{K03.09.minkowski}{Minkowski's finite integral inequality}
\begin{theorem}
For $p\ge1$, write $\|f\|_p=(\int_a^b|f|^p)^{1/p}$.
For integrable scalar $f,g$, $\|f+g\|_p\le\|f\|_p+\|g\|_p$.
\end{theorem}
\begin{proof}
For $p=1$ integrate $|f+g|\le|f|+|g|$.
For $p>1$ put $H=\|f+g\|_p$. All functions below are integrable
by products and continuous composition. If $H=0$ there is nothing to
prove. Otherwise the triangle inequality followed by H\"older with
$q=p/(p-1)$ gives
\begin{align*}
H^p&\le\int(|f|+|g|)|f+g|^{p-1}\\
&\le(\|f\|_p+\|g\|_p)
\left(\int|f+g|^{(p-1)q}\right)^{1/q}
=(\|f\|_p+\|g\|_p)H^{p-1}.
\end{align*}
Divide by $H^{p-1}$. This is a seminorm on actual functions: a function
nonzero at just one point can have zero seminorm. No quotient space is
being constructed here.
\end{proof}

\KEntry{K03.09.partial_fractions}{Rational integration and the algebraic interface}
\begin{construction}
The shared algebra interface used here is the following precise result:
for real polynomials $P,Q$, $Q\ne0$, Euclidean division and factorization
of $Q$ into real linear factors and irreducible real quadratic factors
provide a finite identity
\[\frac P Q=A(x)+\sum_{r,k}\frac{a_{rk}}{(x-r)^k}
+\sum_{j,k}\frac{b_{jk}x+c_{jk}}{(x^2+u_jx+v_j)^k},
\qquad u_j^2-4v_j<0.
\]
This factorization and partial-fraction algebra are accepted from the
fixed shared algebra foundation (D07), not asserted as newly proved D03
theorems. The identity holds where $Q\ne0$.
\end{construction}
\begin{proof}[The analytic reduction]
Polynomial terms integrate by powers. Linear terms integrate to
$a\log|x-r|$ for $k=1$ and
$a(x-r)^{1-k}/(1-k)$ for $k>1$ on intervals excluding $r$.
For a quadratic write $y=x+u/2$, $d=v-u^2/4>0$ and split the numerator
as $By+C$. The $By$ term has primitive
$(B/2)\log(y^2+d)$ if $k=1$, and
$B(y^2+d)^{1-k}/(2(1-k))$ otherwise.
For the constant term define $J_k=\int\dd y/(y^2+d)^k$.
$J_1=d^{-1/2}\arctan(y/\sqrt d)$, and differentiating
$y/(y^2+d)^{k-1}$ gives
\[\left(\frac{y}{(y^2+d)^{k-1}}\right)'
=\frac{2(k-1)d}{(y^2+d)^k}-\frac{2k-3}{(y^2+d)^{k-1}}.
\]
Therefore
\[J_k=\frac{y}{2(k-1)d(y^2+d)^{k-1}}
+\frac{2k-3}{2(k-1)d}J_{k-1}\quad(k\ge2).
\]
This finite recurrence completes the reduction to rational functions,
logarithms and arctangents, and differentiating verifies every term.
Complex partial fractions can instead be reduced on real intervals
componentwise, or by this real decomposition of their real and imaginary
rational parts; a complex logarithm branch is not introduced implicitly.
\end{proof}

\KEntry{K03.09.area_length}{Area and length of a scalar graph}
\begin{theorem}
For continuous $f\ge0$ on $[a,b]$, the area of its subgraph is $\int_a^bf$
when area is defined by inner and outer finite axis-parallel rectangle
contents. For $f\in C^1([a,b])$, its graph has polygonal length
\[L=\int_a^b\sqrt{1+f'(x)^2}\dd x.
\]
\end{theorem}
\begin{proof}
For area, a partition gives interior columns of heights $\inf_J f$ and
covering columns of heights $\sup_J f$. Their total areas are $L(f,P)$
and $U(f,P)$. Boundaries can be covered by finitely many rectangles
of arbitrarily small total area; a common finite grid proves the elementary
inequality that every inner finite rectangle sum is at most every outer
one (decompose all overlaps by their finitely many horizontal and vertical
edges). The inner content is thus bounded above by every $U(f,P)$ and
the outer content below by every $L(f,P)$, while these column constructions
give the reverse bounds in the limit. The Darboux gap tends to zero,
so both contents equal $\int f$.

Define graph length as the supremum over partitions of the sum of chord
lengths for $\gamma(x)=(x,f(x))$. Since
$\gamma(y)-\gamma(x)=\int_x^y(1,f')$, the vector norm estimate gives
every chord length at most $\int_x^y\sqrt{1+f'^2}$. Hence all polygonal
lengths are bounded above by the displayed integral. Conversely, for a
tagged fine partition, uniform continuity of $f'$ gives
\[\left\|\gamma(x_i)-\gamma(x_{i-1})
-(1,f'(\xi_i))\Delta x_i\right\|
\le\omega_{f'}(|P|)\Delta x_i.
\]
The reverse triangle inequality bounds the difference of polygonal length
and the tagged sum for $\sqrt{1+f'^2}$ by
$(b-a)\omega_{f'}(|P|)\to0$. These polygonal lengths approach the
integral, proving the opposite inequality for their supremum.
\end{proof}

\KEntry{K03.09.derivative_not_integrable}{A primitive with nonintegrable derivative}
\begin{example}
Define $F(x)=x^2\sin(1/x^2)$ for $x\ne0$ and $F(0)=0$ on $[-1,1]$.
Then $F'(0)=\lim_{h\to0}h\sin(1/h^2)=0$, while for $x\ne0$,
\[F'(x)=2x\sin(1/x^2)-\frac2x\cos(1/x^2).
\]
At $x_n=(2\pi n)^{-1/2}$ this equals $-2/x_n\to-\infty$.
Thus $F$ is a primitive of its everywhere defined derivative, but the
derivative is unbounded on $[-1,1]$ and is not properly Riemann integrable.
This is the unbounded-derivative example of \KRef{K03.04.unbounded_derivative},
and shows why existence of a primitive does not remove the integrability
condition from Newton--Leibniz.
\end{example}
